\documentclass[a4paper,11pt]{article}

\usepackage[utf8]{inputenc}
\usepackage[T1]{fontenc}
\usepackage{amsmath,amssymb,amsthm}
\usepackage{mathtools}
\usepackage{bm}
\usepackage{hyperref}
\usepackage[margin=1.2in]{geometry}
\usepackage{xcolor}
\usepackage{graphicx}

\newtheorem{proposition}{Proposition}
\newtheorem{lemma}[proposition]{Lemma}
\newtheorem{corollary}[proposition]{Corollary}
\theoremstyle{remark}
\newtheorem{remark}[proposition]{Remark}
\theoremstyle{definition}
\newtheorem{definition}[proposition]{Definition}
\newtheorem{hypothesis}{Hypothesis}

\newcommand{\todo}[1]{\textcolor{red}{[TODO: #1]}}
\newcommand{\planned}[1]{\textcolor{blue}{[PLANNED: #1]}}
\newcommand{\R}{\mathbb{R}}
\newcommand{\E}{\mathbb{E}}
\newcommand{\vect}{\operatorname{vec}}
\newcommand{\tr}{\operatorname{tr}}
\newcommand{\SO}{\mathrm{SO}(3)}
\newcommand{\ip}[2]{\langle #1, #2 \rangle}

\title{The hidden representations and interpretation of SO(3) group operation in neural networks\\[0.5em]
\large Working draft}
\author{\todo{authors}}
\date{\today}

\begin{document}
\maketitle

\begin{abstract}
It is a famous fact that networks trained on finite-group composition with one-hot
inputs learn hidden representations that are matrix elements of single irreducible
representations (irreps). We argue that this fact combines three distinct
ingredients, and use the compact Lie group $\SO$ to separate them.
\emph{Availability}: for a finite group the natural encoding is the regular
representation, suggested by the one-hot vectors, which contains every irrep with
full multiplicity, linearly and up to a change of basis a first layer cannot
detect. For $\SO$ the regular representation $L^2(\SO)$ is infinite-dimensional,
and every encoding built from a finite-dimensional representation is a truncated
inverse Fourier transform; we use the truncated delta
$\bigoplus_{l\le L}\sqrt{2l+1}\,\vect D^l(G)$, also presented in a ``pixel'' basis of
samples at random anchor rotations, where no irrep is visible. \emph{Target structure}: composition is block-diagonal in
Fourier space, $D^l(GH)=D^l(G)D^l(H)$. \emph{Selection}: an elementwise
nonlinearity must choose a basis of hidden units; we find that irrep-pure units are selected whenever the norm is controlled,
explicitly by weight decay or implicitly by a small initialisation, while explicit
weight decay is needed to prune them to the minimum-norm count. With a
quadratic activation we show that (i) when an encoding exposes an irrep only
through a rank-$r$ Fourier block, a hidden layer can capture at most a fraction
$r/(2l+1)$ of the corresponding target, a bound reached by trained networks;
(ii) with the truncated delta as input and output, every active neuron is
irrep-pure, with the same irrep in both input blocks and the readout, in both the
Fourier and the pixel basis; (iii) within each irrep the network implements
$d_l\times d_l$ matrix multiplication (tensor cosine $0.96$--$1.000$), with a
causal, block-diagonal ablation profile; and (iv) under weight decay it converges
to a sharply defined number of equal-norm rank-one product units
$(a^\top D^l(G)u)(u^\top D^l(H)b)\,ab^\top$ per irrep: $d_l^3$ for $l=1,2$ and
$d_l^3+d_l$ for $l=3,4$. The irreps are thus not discovered from the encoding;
they are selected among the bases an elementwise network could use.
\end{abstract}

\tableofcontents

%======================================================================
\section{Introduction}
%======================================================================

Neural networks trained to compose elements of a finite group $G$ learn
features built from the group's irreducible representations: for modular
addition, neurons are Fourier modes \cite{nanda2023progress, gromov2023grokking,
zhong2023clock}; for general finite groups, neurons are real linear
combinations of the matrix elements of a single irrep \cite{chughtai2023toy,
stander2024cosets}. For two-layer networks with polynomial activation and small
initialisation, Marchetti et al.\ \cite{marchetti2026sequential} prove this,
using the Alternating Gradient Flow (AGF) framework \cite{kunin2025agf}: irreps
are acquired one at a time, in an order set by the Fourier statistics of the
input encoding. Margin-maximisation arguments give a related prediction for
quadratic networks \cite{morwani2023feature}.

The phenomenon is usually told as ``the network discovers the Fourier transform
of the group''. We argue that this conflates three ingredients:
\begin{enumerate}
\item \textbf{Availability.} For a finite group, the natural encoding is the
  regular representation $\R[G]$, suggested by the one-hot vectors, in which every
  irrep $\rho$ occurs with multiplicity $d_\rho$; its Fourier transform, an
  isometry, maps the encoding of $g$ to the matrix coefficients
  $\sqrt{d_\rho/|G|}\,\rho(g)$ (Section~\ref{sec:finite}). A first layer with
  isotropic initialisation and weight decay is blind to that change of basis, so
  the irreps are linearly available from the start.
\item \textbf{Target structure.} Composition is a convolution, hence
  block-diagonal in Fourier space: $\rho(gh)=\rho(g)\rho(h)$ for every irrep.
\item \textbf{Selection.} With an elementwise activation, rotating the hidden
  units changes the function, so the network must pick a basis. That it picks
  units that are pure in one irrep, in both inputs and in the readout, is the
  genuinely non-trivial part.
\end{enumerate}

Finite groups make these ingredients hard to pull apart, because the regular
representation is both natural and finite-dimensional. For the compact Lie group
$\SO$ it is $L^2(\SO)$, infinite-dimensional, so it cannot serve as an encoding:
any encoding built from a finite-dimensional representation is a truncated inverse
Fourier transform, involving finitely many irreps with some coefficients. The
choice of encoding, rotation matrices, orbits $\rho(G)x$ of a fixed vector, or the
delta function at $G$ truncated to irreps $l\le L$, then decides which irreps are
available, and $\SO$ becomes a testbed where the three ingredients separate. The truncated delta can be
given in two bases. In the \emph{Fourier basis} it is the family of its Peter--Weyl
coefficients, where the irreps are explicit. In the \emph{pixel basis} it is the
vector of its values at $N$ fixed anchor rotations $R_1,\dots,R_N$, a blurred bump
centred at $G$; this is the analogue of the one-hot coordinates of a finite group,
which are the values of the delta function at each group element, and no irrep is
visible in it. Finally, the target can contain one irrep (the product $GH$) or many
(the truncated delta of $GH$).

\paragraph{Contributions.}
\begin{itemize}
\item \textbf{Encodings} (Section~\ref{sec:encodings}). For a finite group the
  regular representation is a natural, finite-dimensional encoding; for $\SO$ it is
  infinite-dimensional, and every encoding built from a finite-dimensional
  representation is a truncated inverse Fourier transform. We use the truncation
  closest to the regular representation, the truncated delta
  $\delta^L(G)=\bigoplus_{l\le L}\sqrt{2l+1}\vect D^l(G)$ as a truncated Fourier
  transform, its pixel-basis version (samples at Haar-random anchors), and
  characterise its convergence (Figure~\ref{fig:delta}). Rotation matrices are
  related to it by the degree filtration: $D^l$ is a polynomial of degree exactly
  $l$ in the entries of $G$.
\item \textbf{Availability bound} (Section~\ref{sec:capacity}). Consider an
  orbit encoding $u(G)=\rho(G)x$ of a fixed vector $x$ in a representation $\rho$
  of $\SO$. The $l$-isotypic component of $x$ is a $d_l\times m_l$ matrix $X_l$
  ($m_l$ the multiplicity of $l$ in $\rho$), and the input contains $D^l(G)$ only
  through the product $D^l(G)X_l$, i.e.\ only $r_l=\operatorname{rank}X_l$ of the
  $d_l=2l+1$ directions on which $D^l(G)$ acts. We prove that a network with one
  quadratic hidden layer, of any width, then explains at most a fraction
  $r_l/d_l$ of the variance of the $l$-component of the target $\rho(GH)x$
  (Proposition~\ref{prop:capacity}). There is no loss when every $X_l$ has full
  rank, which holds for the regular representation (and hence for the one-hot
  encoding of a finite group), but not for tensor powers of $\R^3$, where
  $m_l<d_l$ for large $l$. For the matrix target $GH$ the bound becomes
  $(r_1/d_1)^2$, and trained networks attain it: with $r_1=1$ they explain $0.113$ of
  the variance, against the bound $1/9\approx0.111$.
\item \textbf{Mechanism} (Sections~\ref{sec:mechanism}--\ref{sec:experiments}).
  Schur orthogonality predicts the per-neuron signature of the quadratic product
  mechanism (Lemma~\ref{lem:cross}); trained networks follow it, computing
  $D^l(GH)$ as a sum of rank-one products of $D^l(G)$ and $D^l(H)$.
\item \textbf{Selection} (Section~\ref{sec:planned}). With the truncated delta as
  input and output, all active neurons are irrep-pure and consistent, in the
  Fourier and in the pixel basis; under weight decay each irrep is served by
  $d_l^3$ ($l\le2$) or $d_l^3+d_l$ ($l=3,4$) equal-norm rank-one product units.
\end{itemize}

%======================================================================
\section{Background}
%======================================================================

\paragraph{Notation.} $G,H\in\SO$ are sampled from the Haar measure. For each
$l\ge0$, the irrep of degree $l$ of $\SO$ acts on the $d_l=2l+1$ dimensional space
of spherical harmonics of degree $l$, by $(G\cdot f)(x)=f(G^{-1}x)$. Its matrix in
the orthonormal basis of real spherical harmonics $Y_l^{-l},\dots,Y_l^{l}$ is the
\emph{real Wigner matrix} $D^l(G)\in O(d_l)$, defined by
$Y_l^m(G^{-1}x)=\sum_{m'}D^l_{m'm}(G)\,Y_l^{m'}(x)$; $G\mapsto D^l(G)$ is a
homomorphism, all irreps of $\SO$ are of real type, and $D^1(G)=G$ up to a fixed
permutation of axes. Explicit formulas and the computation we use are in
Appendix~\ref{app:wigner}. We write $\ip{A}{M}=\tr(A^\top M)=\sum_{ik}A_{ik}M_{ik}$.

\paragraph{Schur orthogonality.} Under Haar measure,
\begin{equation}\label{eq:schur}
\E_G\bigl[D^l_{ij}(G)\,D^{l'}_{mn}(G)\bigr]=\frac{\delta_{ll'}\delta_{im}\delta_{jn}}{d_l}.
\end{equation}

\paragraph{Peter--Weyl.} $L^2(\SO)=\bigoplus_{l\ge0}\mathcal{H}_l$, where
$\mathcal{H}_l$ is spanned by the $d_l^2$ functions $G\mapsto D^l_{mn}(G)$;
as a left representation, $\mathcal{H}_l$ contains the irrep $l$ with
multiplicity $d_l$. The same holds for a finite group, with $L^2(G)=\R[G]$
and the multiplicity of $\rho$ equal to $d_\rho$.

\paragraph{Casimir and Laplacian.} The Casimir operator of $\SO$ acts on
$\mathcal{H}_l$ as the scalar $l(l+1)$; it coincides with the
Laplace--Beltrami operator of the bi-invariant metric. Hence the $D^l_{mn}$ are
Laplacian eigenfunctions, and isotypic projectors in any representation can be
computed exactly from the Casimir (we use this throughout).

%======================================================================
\section{Encodings: from the regular representation to the truncated delta}
\label{sec:encodings}
%======================================================================

A group element $g$ has a canonical representative, its image under the regular
representation: the delta function $\delta_g=\lambda(g)\delta_e$ on $G$. For a
finite group $\R[G]$ is finite-dimensional and $\delta_g$ is the one-hot vector, so
the regular representation is the natural encoding, suggested by the one-hot
vectors themselves. For a compact Lie group the regular representation is
$L^2(G)$, which is infinite-dimensional, and $\delta_g$ is not even an element of it:
it cannot be given to a network as such. This is the essential difference between
the two cases, and it shapes the rest of this section.

\subsection{Finite groups}\label{sec:finite}

Let $G$ be finite with real irreps $\rho$ (for simplicity of real type;
complex and quaternionic types require the usual realification). Define the
Fourier transform $F:\R^{|G|}\to\bigoplus_\rho\R^{d_\rho\times d_\rho}$,
\begin{equation}
(Fv)_\rho=\sqrt{\tfrac{d_\rho}{|G|}}\sum_{h\in G}v(h)\,\rho(h).
\end{equation}

By the Plancherel theorem $F$ is orthogonal. The one-hot vector of $g$ is
$e_g=\lambda(g)e_1$, where $\lambda$ is the left-regular representation on
$\R[G]\cong\R^{|G|}$, and
\begin{equation}\label{eq:onehot}
Fe_g=\Bigl(\sqrt{\tfrac{d_\rho}{|G|}}\;\rho(g)\Bigr)_{\rho},\qquad
F(\lambda(g)x)_\rho=\rho(g)\,\hat x(\rho),\quad \hat x(\rho)=(Fx)_\rho .
\end{equation}

\begin{remark}[Consequences for networks]\label{rem:onehot}
A first layer $W e_g=(WF^\top)(Fe_g)$ is the same network written in the
Fourier basis; since $F$ is orthogonal, the Frobenius norm of $W$, and hence
weight decay and isotropic initialisation, are unchanged. For such a layer, a
one-hot input is therefore equivalent to the family of irrep matrix coefficients of
$g$. This removes much of the surprise usually attached to irrep features: the
network is discovering a Fourier transform but, the irrep
matrix coefficients are only a change of basis of the input. Every linear function of the one-hot input is already a
combination of irrep matrix elements. What remains to be explained is why each
hidden unit uses the coefficients of a \emph{single} irrep rather than a mixture
(Section~\ref{sec:planned}). For an orbit encoding $\lambda(g)x$
the input is $\rho(g)\hat x(\rho)$: whenever every $\hat x(\rho)$ is
invertible, $\rho(g)$ is a linear function of the input, and the encoding is
equivalent to one-hot up to an (in general non-orthogonal) linear map. Since
$\R[G]$ contains $\rho$ with multiplicity $d_\rho$, $\hat x(\rho)$ is a
\emph{square} matrix and is invertible for generic $x$. This is why random
orbit encodings work as well as one-hot for finite groups
\cite{marchetti2026sequential}.
\end{remark}

\subsection{$\SO$: every encoding is a truncated inverse Fourier transform}

Any encoding of $\SO$ built from a finite-dimensional representation, that is an
orbit map $G\mapsto\rho(G)x$ with $\dim\rho<\infty$ (rotation matrices, tensor
powers acting on a fixed tensor, the truncated delta below), only involves the
finitely many irreps contained in $\rho$. Decomposing $\rho=\bigoplus_l
I_{m_l}\otimes D^l$, the encoding is, up to a fixed linear map,
\[
\rho(G)x\;\cong\;\bigoplus_{l\in S}D^l(G)\,X_l ,
\]
a truncated inverse Fourier transform: a finite set $S$ of irreps, each seen
through the coefficient matrix $X_l$ of $x$. Unlike the finite case, truncation is
therefore not a choice but a necessity; what can be chosen is which irreps enter,
and with which coefficients. Non-linear charts (Euler angles, exponential
coordinates) escape this description but are not representations, and the group
action on them is not linear.

The Peter--Weyl expansion of the delta function at $G$ is
\begin{equation}
\delta_G(H)=\sum_{l\ge0}d_l\,\tr\bigl(D^l(G)^\top D^l(H)\bigr),
\end{equation}
i.e.\ the coefficients of $\delta_G$ in the Peter--Weyl basis are
$\bigl(\sqrt{d_l}\,D^l(G)\bigr)_{l\ge0}$. Truncating this expansion at degree $L$ is
the choice closest to the regular representation: it keeps every irrep up to $L$,
each with full-rank coefficients ($X_l=\sqrt{d_l}\,I$) and the same weights as the
delta itself.

\begin{definition}[Truncated delta]
For $L\ge1$, $\delta^L(G)=\bigoplus_{l=1}^{L}\sqrt{d_l}\,\vect D^l(G)\in
\R^{\sum_{l=1}^L d_l^2}$ (the constant $l=0$ block is dropped).
\end{definition}

Its Gram kernel $\ip{\delta^L(G)}{\delta^L(H)}=\sum_{l\le L}d_l\,\chi_l(G^\top
H)$ is a Dirichlet-type kernel concentrated at $G=H$: a blurred delta.
With the $\sqrt{d_l}$ weights, a first-layer weight block $W^{(l)}$ produces
the preactivation $\sqrt{d_l}\ip{W^{(l)}}{D^l(G)}$, whose Haar variance is
exactly $\|W^{(l)}\|_F^2$ by \eqref{eq:schur}: weight norm per block equals
function energy per irrep.

The truncated delta can be presented in two bases. In the Fourier basis above the
irreps are explicit. Sampling the function $\delta^L_G=\sum_{l\le L}d_l\chi_l(G^\top\cdot)$
at $N$ fixed anchor rotations instead gives the \emph{pixel basis}, a blurred
one-hot over the anchors in which no irrep is visible; this is the $\SO$
counterpart of the one-hot vectors of a finite group, and it is the encoding that
justifies, through the truncated delta, the use of a ``pixelised one-hot''
(Section~\ref{sec:computing-encoding}).

\subsection{Computing the encoding: a truncated Fourier transform of the delta}
\label{sec:computing-encoding}

The encoding is the Fourier transform on $\SO$ of the delta at $G$, truncated at
degree $L$. For $f\in L^2(\SO)$ the Fourier coefficient at degree $l$ is the
$d_l\times d_l$ matrix
\[
\hat f(l)=\sqrt{d_l}\int_{\SO} f(R)\,D^l(R)\,dR ,
\]
and with this normalisation $f\mapsto(\hat f(l))_{l\ge0}$ is an isometry
(Plancherel). Applied to $f=\delta_G$, the integral collapses to
$\hat\delta_G(l)=\sqrt{d_l}\,D^l(G)$: \emph{the Fourier transform of a delta is
the irrep matrix itself}, so no integral has to be computed. Encoding a rotation
therefore amounts to:
\begin{enumerate}
\item \textbf{Real Wigner matrices.} For $l=1,\dots,L$, compute the real
  orthogonal matrix $D^l(G)\in O(d_l)$. Any standard method works; three
  equivalent ones:
  \begin{itemize}
  \item \emph{via spherical harmonics}: $D^l(G)$ is the matrix of
    $Y\mapsto Y\circ G^{-1}$ on the real harmonics of degree $l$; evaluate the
    $Y_l^{m}$ on a quadrature grid of $S^2$ and on the rotated grid, and solve for
    the coefficients (our implementation);
  \item \emph{via Euler angles}: $D^l(\alpha,\beta,\gamma)$ from the closed form
    $e^{-im\alpha}d^l_{mn}(\beta)e^{-in\gamma}$, then the fixed unitary change to
    the real basis;
  \item \emph{via tensor powers}: $D^l(G)=P_l^\top G^{\otimes l}P_l$, where the
    columns of $P_l$ are an orthonormal basis of the symmetric traceless tensors
    of order $l$ (the harmonic polynomials of degree $l$). This makes explicit
    that the entries of $D^l$ are polynomials of degree exactly $l$ in the
    entries of $G$.
  \end{itemize}
\item \textbf{Scale and concatenate.}
  $\delta^L(G)=\bigl(\sqrt{d_l}\,\vect D^l(G)\bigr)_{l=1}^{L}\in\R^{N_L}$, with
  $N_L=\sum_{l=1}^Ld_l^2=\tfrac{(L+1)(2L+1)(2L+3)}{3}-1$ ($83$ for $L=3$, $164$
  for $L=4$). The $l=0$ block is the constant $1$ and is dropped.
\item \textbf{Target.} For composition, the target is $\delta^L(GH)$; blockwise
  $D^l(GH)=D^l(G)D^l(H)$, so it can be computed either from the product $GH$ or
  from the two encodings.
\end{enumerate}

\paragraph{Properties.} (i) Every entry of $\delta^L(G)$ has unit Haar variance
and zero mean (Schur), and $\|\delta^L(G)\|^2=N_L$ for every $G$. (ii) A
first-layer weight block $W^{(l)}$ gives a preactivation with Haar variance
exactly $\|W^{(l)}\|_F^2$, which makes per-neuron irrep spectra fit-free.
(iii) The Gram kernel is the truncated delta:
$\langle\delta^L(G),\delta^L(H)\rangle=K_L(G,H)-1$ with
\[
K_L(G,H)=\sum_{l=0}^{L}d_l\,\chi_l(\theta),\qquad
\chi_l(\theta)=\frac{\sin\bigl((2l+1)\theta/2\bigr)}{\sin(\theta/2)},
\]
$\theta$ the angle of $G^\top H$. In closed form, with $x=\theta/2$,
$K_L=-\frac{1}{\sin x}\frac{d}{dx}\Bigl[\frac{\sin(2(L+1)x)}{2\sin x}\Bigr]$, the
$\SO$ analogue of the Dirichlet kernel. It has unit mass for every $L$
($\int\chi_l=\delta_{l0}$), height $K_L(0)=\sum_{l\le L}d_l^2\sim\frac43L^3$, and
width $\sim1/L$; since Haar measure near the identity is $\propto\theta^2d\theta$,
it converges to $\delta_G$ as $L\to\infty$ in the sense of distributions: by
Peter--Weyl completeness, $\int f\,K_L(\cdot,H)\to f(H)$ for smooth $f$
(Figure~\ref{fig:delta}a). The convergence is genuinely weaker than for a
positive approximate identity: $K_L$ is not positive, and its tail, weighted by
the Haar density $\propto\theta^2$, keeps an $O(1)$ oscillation, so the mass
within angle $\theta$ of the identity does not tend to a step
(Figure~\ref{fig:delta}b). This is the Pinsky phenomenon of Fourier series in
dimension $\ge3$. Weighting the irreps by $e^{-t\,l(l+1)}$ (the heat kernel on
$\SO$, $l(l+1)$ being the Laplacian eigenvalue) gives a positive kernel whose mass
does concentrate at the identity (Figure~\ref{fig:delta}c).

\begin{figure}[h]
\centering
\includegraphics[width=\textwidth]{figures/delta_kernel.pdf}
\caption{The truncated delta on $\SO$ as a function of the rotation angle
$\theta$ of $G^\top H$, for $L=1,2,4,8,16$ (light to dark). (a)
$K_L(\theta)/K_L(0)$: the peak narrows, with width $\sim1/L$ and height
$K_L(0)=\sum_{l\le L}d_l^2\sim\frac43L^3$. (b) Haar mass within angle $\theta$
for the sharp truncation: it reaches 1 at $\theta=\pi$ for every $L$ but keeps
oscillating, so $K_L\to\delta$ only against smooth test functions. (c) The same
with heat-kernel weights $e^{-t\,l(l+1)}$, $t=8/L^2$: the kernel is positive and
its mass concentrates at the identity. Script:
\texttt{scripts/plot\_delta\_kernel.py}.}
\label{fig:delta}
\end{figure}

\paragraph{Position-basis version.} Sampling the truncated delta at fixed anchor
rotations $R_1,\dots,R_N$, $x_G[i]=K_L(G,R_i)$, gives an input that looks like a
blurred one-hot over the anchors. If the anchors form a quadrature exact for
degree $2L$, this is an orthogonal image (up to scale) of $\delta^L(G)$ and hence
an equivalent input for a first layer with isotropic initialisation and weight
decay, exactly as one-hot and Fourier coordinates are equivalent for finite
groups (Remark~\ref{rem:onehot}).

\paragraph{Tensor powers.} Decompose $(\R^3)^{\otimes k}=\bigoplus_l
\R^{m_l}\otimes\R^{d_l}$, so that for an orthogonal $U$,
\begin{equation}\label{eq:tensor}
G^{\otimes k}=U\Bigl(\bigoplus_l I_{m_l}\otimes D^l(G)\Bigr)U^\top,
\end{equation}
with $(m_0,m_1)=(0,1)$ for $k=1$, $(m_0,m_1,m_2)=(1,1,1)$ for $k=2$, and
$(m_0,\dots,m_3)=(1,3,2,1)$ for $k=3$.
\begin{itemize}
\item \emph{Matrix encoding} $\vect(G^{\otimes k})$: an orthogonal image of
  $\bigl(D^l(G)\bigr)_{l\le k}$, each block repeated $m_l$ times. It is
  $\delta^k(G)$ up to a linear map and block weights ($\sqrt{m_l}$ instead of
  $\sqrt{d_l}$), which change initialisation and weight-decay geometry but not
  expressivity.
\item \emph{Orbit encoding} $G^{\otimes k}x$ with a fixed $x$: its $l$-block is
  $D^l(G)\,\hat x_l$ with $\hat x_l\in\R^{d_l\times m_l}$, of rank
  $r_l\le\min(m_l,d_l)$. For $k=3$: $r_1=3=d_1$ but $r_2\le2<5$, $r_3\le1<7$.
  This is the analogue of a finite-group orbit encoding with rank-deficient
  Fourier blocks, which never happens in $\R[G]$ but always happens here,
  because $(\R^3)^{\otimes k}$ contains too few copies of the high irreps.
\end{itemize}

\begin{remark}[How $\ip{A}{G}$ is read from $G^{\otimes 3}x$]
Contracting two indices commutes with rotation: $\sum_b (G^{\otimes3}x)_{abb}
=(Gt_1)_a$ with $t_1=\sum_b x_{\cdot bb}$, and similarly $Gt_2,Gt_3$ for the
other index pairs. A weight vector $w_{abc}=a_a\delta_{bc}$ thus computes
$a^\top Gt_1=\ip{at_1^\top}{G}$, and for generic $x$ the $t_m$ are independent,
so every $\ip{A}{G}$ is linear in the input. For $k=2$ the only one-index
contraction is the axial vector $c$ of the antisymmetric part
($G[c]_\times G^\top=[Gc]_\times$), so only $\ip{ac^\top}{G}$ is available; for
$k=1$ only $\ip{ax^\top}{G}$, and $G$ is not even identifiable from $Gx$.
\end{remark}

\begin{proposition}[Degree filtration]\label{prop:degree}
The entries of $D^l(G)$ are polynomials of degree exactly $l$ in the entries of
$G$, and the polynomials of degree $\le k$ in the entries of $G$, restricted to
$\SO$, span exactly the matrix elements of the irreps $l\le k$.
\end{proposition}

\begin{proof}[Proof sketch]
The irrep $l$ occurs in $(\R^3)^{\otimes l}$ (as the symmetric traceless tensors),
so $D^l(G)=P_l^\top G^{\otimes l}P_l$ for a fixed isometry $P_l$: degree $\le l$.
Conversely, a polynomial of degree $k$ in the entries of $G$ is a matrix element of
$\bigoplus_{j\le k}(\R^3)^{\otimes j}$, which only contains irreps $l\le k$
(Clebsch--Gordan), so by Schur orthogonality it has no component on $D^l$ for
$l>k$; hence the degree of $D^l$ is exactly $l$. More generally, for any compact
group, the matrix elements of a faithful representation and its dual generate all
representative functions (Peter--Weyl and Stone--Weierstrass); for $\SO$ the
faithful defining representation coincides with the adjoint one.
\end{proof}

Numerically, fitting the entries of $D^l(G)$ by polynomials of degree $\le k$ in the
entries of $G$ gives held-out $R^2=0$ for $k<l$ and $R^2=1.000$ for $k\ge l$
($l\le4$). Rotation matrices therefore expose only $l=1$ linearly: higher irreps
must be \emph{computed}, as degree-$l$ products, while the truncated delta makes
them all linearly available. A network of depth $D$ with quadratic activations is a
polynomial of degree $2^D$ in its inputs, so with rotation-matrix inputs it can
represent the $l$-block of $\delta^L(GH)$ (degree $l$ in $G$ and in $H$) only if
$2l\le2^D$.

%======================================================================
\section{An availability bound for quadratic layers}
\label{sec:capacity}
%======================================================================

Consider one hidden layer with $\sigma(z)=z^2$, no normalisation, arbitrary
width:
\begin{equation}\label{eq:net}
f(G,H)=\sum_j r_j\bigl(\ip{w_j}{u(G)}+\ip{v_j}{u(H)}+b_j\bigr)^2+c ,
\end{equation}
where $u(\cdot)$ is the input encoding. Expanding, $f$ is a constant plus terms
depending on $G$ only, terms depending on $H$ only, and the cross part
$2\sum_j r_j\ip{w_j}{u(G)}\ip{v_j}{u(H)}$, an arbitrary bilinear form in
$(u(G),u(H))$ when the width is large enough.

\begin{proposition}[Capacity of a quadratic layer]\label{prop:capacity}
Let $u(G)=\rho(G)x$ and the target be the orbit $y(G,H)=\rho(GH)x$, with
$\rho=(\R^3)^{\otimes k}$. For $l\ge1$, the fraction of the Haar variance of
the $l$-component of $y$ that any $f$ of the form \eqref{eq:net} can capture is
at most $r_l/d_l$, $r_l=\operatorname{rank}\hat x_l$, and this is attained.
For the matrix target $y=GH$ ($l=1$) with an orbit input, the bound is
$(r_1/d_1)^2$.
\end{proposition}

\begin{proof}
Write the $l$-isotypic part of $\rho$ as $I_{m_l}\otimes D^l$ and collect the $m_l$
copies of $x$'s $l$-component as the columns of $X_l\in\R^{d_l\times m_l}$, so that
the $l$-part of the input is the matrix $D^l(G)X_l$ and the $l$-part of the target
is $Y_l=D^l(G)D^l(H)X_l$; its Haar second moment is $\E\|Y_l\|_F^2=\|X_l\|_F^2$.

\emph{Reduction to bilinear features.} Expanding~\eqref{eq:net}, $f$ is a constant,
plus functions of $G$ alone and of $H$ alone, plus a bilinear form in
$(u(G),u(H))$; conversely every bilinear form is realised by finitely many units,
since $(a+b)^2-(a-b)^2=4ab$ removes the non-bilinear terms. For $l\ge1$,
$\E_G[D^l(G)]=0$, so functions of $G$ alone (or $H$ alone) and constants are
orthogonal to $Y_l$. By~\eqref{eq:schur}, products of entries of $D^{l'}(G)X_{l'}$
and $D^{l''}(H)X_{l''}$ with $(l',l'')\ne(l,l)$ are orthogonal to $Y_l$ as well. It
therefore suffices to project $Y_l$ onto
$\mathcal V_l=\operatorname{span}\{(D^l(G)X_l)_{ip}(D^l(H)X_l)_{jq}\}$.

\emph{The projection.} Let $\Pi_l=X_lX_l^+$ be the orthogonal projector onto the
column space of $X_l$, of rank $r_l$. The candidate
$\hat Y_l=D^l(G)\,\Pi_l\,D^l(H)X_l$ lies in $\mathcal V_l$, since
$(\hat Y_l)_{iq}=\sum_{p,j}(D^l(G)X_l)_{ip}(X_l^+)_{pj}(D^l(H)X_l)_{jq}$. The
residual $Y_l-\hat Y_l=D^l(G)(I-\Pi_l)D^l(H)X_l$ is orthogonal to $\mathcal V_l$:
integrating over $G$ with~\eqref{eq:schur},
\[
\E_G\bigl[(D^l(G)(I-\Pi_l)D^l(H)X_l)_{iq}\,(D^l(G)X_l)_{i'p}\bigr]
=\tfrac{\delta_{ii'}}{d_l}\bigl(X_l^\top(I-\Pi_l)D^l(H)X_l\bigr)_{pq}=0,
\]
because $X_l^\top(I-\Pi_l)=0$. Hence $\hat Y_l$ is the best approximation, and
\[
\E\|\hat Y_l\|_F^2=\E\|\Pi_lD^l(H)X_l\|_F^2
=\sum_q\E\bigl\|\Pi_lD^l(H)x_q\bigr\|^2=\frac{\tr\Pi_l}{d_l}\sum_q\|x_q\|^2
=\frac{r_l}{d_l}\,\|X_l\|_F^2,
\]
using $\E[D^l(H)vv^\top D^l(H)^\top]=\|v\|^2I/d_l$. The captured fraction is
$r_l/d_l$, attained by the best bilinear form. For the matrix target $GH$ the
same computation applies on both sides: the $H$-features are read through the
column space of $X_1$ on the left of $D^1(H)$ as well, which gives
$\E\|D^1(G)\Pi_1D^1(H)\Pi_1\|_F^2/\E\|D^1(G)D^1(H)\|_F^2=r_1^2/d_1^2$.
\end{proof}

\begin{corollary}
With the truncated delta (or the matrix encoding), $r_l=d_l$ for all
$l\le L$ and the target is exactly representable by a quadratic layer
(e.g.\ by the explicit construction in Section~\ref{sec:mechanism}).
\end{corollary}

\paragraph{Numerical check.} For $k=3$, a random unit $x$, and the orbit target,
we fit the best bilinear-plus-constant function by least squares on $30{,}000$
Haar pairs and evaluate $R^2$ per output irrep on $10{,}000$ held-out pairs:

\begin{center}
\begin{tabular}{lccc}
\hline
& $l=1$ & $l=2$ & $l=3$\\
\hline
measured (1 probe) & 1.00 & 0.385 & 0.124\\
predicted $r_l/d_l$ & 1 & 0.400 & 0.143\\
measured (3 probes) & 1.00 & 1.00 & 0.286\\
predicted $\min(1,3m_l/d_l)$ & 1 & 1 & 0.429\\
\hline
\end{tabular}
\end{center}
The small shortfall is finite-sample shrinkage of the held-out fit (it grows
with the number of features: $730$ for one probe, $6562$ for three).
In trained networks (Section~\ref{sec:experiments}), $k=2$ with the matrix target
reaches relative MSE $0.887$ on the bilinear part, against the bound
$1-(1/3)^2=0.889$.

%======================================================================
\section{The quadratic-product mechanism}
\label{sec:mechanism}
%======================================================================

\subsection{An explicit construction}

Since $(GH)_{ij}=\sum_kG_{ik}H_{kj}$,
\begin{equation}\label{eq:matmul}
GH=\sum_{i,j,k}E_{ij}\,\ip{E_{ik}}{G}\ip{E_{kj}}{H},
\end{equation}
where $E_{ik}$ are matrix units: 27 ``product units'' $(A,B,r)=(E_{ik},E_{kj},E_{ij})$,
each satisfying $r=AB$. A quadratic neuron computes
$(\ip{A}{G}+\ip{B}{H}+c)^2$, whose only useful part is the cross term
$2\ip{A}{G}\ip{B}{H}$; the rest (``extraneous'' terms, cf.\
\cite{marchetti2026sequential}) depends on $G$ or $H$ alone and must cancel.
The polarisation identity $(a+b)^2-(a-b)^2=4ab$ realises each product with two
neurons, so 54 quadratic neurons compute $GH$ exactly. More generally, any
orthonormal basis $\{u_k\}$ gives $GH=\sum_k(Gu_k)(u_k^\top H)$: neurons reading
$a^\top Gu$ and $u^\top Hb$ \emph{with the same inner direction $u$} produce
the useful products.

\subsection{The signature of the mechanism}

\begin{lemma}[Correlation of a cross term with the target]\label{lem:cross}
For $A,B\in\R^{d_l\times d_l}$ and Haar $G,H$,
\[
\E\bigl[D^l(G)D^l(H)x\;\ip{A}{D^l(G)}\ip{B}{D^l(H)}\bigr]=\frac{ABx}{d_l^2},
\qquad
\E\bigl[GH\;\ip{A}{G}\ip{B}{H}\bigr]=\frac{AB}{9}.
\]
\end{lemma}
\begin{proof}
By \eqref{eq:schur}, $\E[D_{ik}(G)\ip{A}{D(G)}]=A_{ik}/d$. Then
$\E[(D(G)D(H)x)_i\ip{A}{D(G)}\ip{B}{D(H)}]=\sum_{k,m}\E[D_{ik}(G)\ip{A}{D(G)}]\,
\E[D_{km}(H)\ip{B}{D(H)}]\,x_m=(ABx)_i/d^2$; the matrix case is identical.
\end{proof}

Since gradient descent on the readout aligns each neuron's readout with the
target's correlation with its feature, the Lemma predicts, for neuron $j$ with
l=1 maps $A_j,B_j$ and effective cross-term weight
$r_j=W^{\text{out}}_{:,j}\,\sigma''(c_j)/2$ ($=W^{\text{out}}_{:,j}$ for $\sigma(z)=z^2$;
a Taylor expansion for smooth activations):
\begin{itemize}
\item[(P1)] \textbf{readout alignment}: $r_j\parallel\vect(A_jB_j)$;
\item[(P2)] \textbf{population}: $T=\sum_jr_j\otimes\vect A_j\otimes\vect B_j$ is
  proportional to the matrix-multiplication tensor $M$, $M_{(ij),(ik),(kj)}=1$;
\item[(P3)] \textbf{composability}: $\|A_jB_j\|_F/(\|A_j\|_F\|B_j\|_F)$ is large
  (inner directions agree);
\item[(P4)] \textbf{cancellation}: the extraneous terms vanish at the
  population level.
\end{itemize}
This is the $\SO$ instance of the finite-group mechanism: preactivations are
irrep matrix elements of each argument, the activation produces their product,
and the readout uses the homomorphism $D^l(G)D^l(H)=D^l(GH)$, as in the
Fourier circuit $\cos(\omega a)\cos(\omega b)\to\cos(\omega(a+b))$ for modular
addition \cite{nanda2023progress}.

%======================================================================
\section{Mechanism: the defining irrep and the target $GH$}
\label{sec:experiments}
%======================================================================

\paragraph{Setup.} Input: orbit encoding $u(G)=G^{\otimes k}x$ with a fixed
unit $x$ ($k\in\{2,3\}$); target $GH\in\R^9$; MSE loss; one hidden layer of
width 2048; activation $z^2$ (SiLU for comparison); AdamW (lr $3\cdot10^{-4}$,
or $10^{-4}$ for SiLU), weight decay $10^{-4}$, batch 128, 1001 epochs,
$10^4$ Haar triples ($80/10/10$ split). Two seeds ($42$, $43$; the seed also
changes $x$) for the main $k=3$ configurations, one otherwise. All analysis
uses fresh Haar samples.

\paragraph{Measurements.}
(i) Geodesic error after SVD projection to $\SO$. (ii) Exact irrep shares of
the first-layer $G$-block preactivation at $H=I$, at initialisation and after
training, via Casimir projectors (no Wigner regression). (iii) Least-squares
fit $z_j\approx\ip{A_j}{G}+\ip{B_j}{H}+c_j$. (iv) P1--P3, each with a
mismatched-neuron null. (v) Output sectors: on a $150\times150$ grid of $(G,H)$
pairs, the output is split into pure-$G$, pure-$H$ and bilinear
(double-centred) parts; the bilinear part is compared with that of $GH$.
(vi) Irrep ablation: remove the $l$-components of $x$ (equivalent to projecting
the first-layer weights).

\begin{table}[h]
\centering\small
\resizebox{\textwidth}{!}{%
\begin{tabular}{lcccccc}
\hline
run & err ($^\circ$) & $l{=}1/2/3$ share init $\to$ final & $l{=}1$ $R^2$ & P3 (null) & P1 (null) & P2 $\cos(T,M)$\\
\hline
$k{=}3$, $z^2$, s42 & \textbf{0.14} & .28/.32/.40 $\to$ .80/.10/.10 & .80 & .89 (.53) & .96 ($-$.01) & 1.000\\
$k{=}3$, $z^2$, s43 & 0.21 & .60/.25/.15 $\to$ .88/.07/.05 & .88 & .75 (.49) & .69 (.00) & 1.000\\
$k{=}3$, SiLU, s42 & 0.23 & .28/.32/.40 $\to$ .85/.07/.08 & .85 & .91 (.53) & .91 ($-$.01) & 1.000\\
\hline
$k{=}2$, $z^2$ & 109.1 & .28/.72 $\to$ .57/.43 & .57 & .78 (.78) & .81 (.01) & 0.332\\
$k{=}2$, SiLU & 109.1 & .28/.72 $\to$ .43/.57 & .43 & .72 (.72) & .84 (.01) & 0.332\\
\hline
\end{tabular}}
\caption{Matrix target $GH$. Chance geodesic error is $126.5^\circ$. Shares are
of the first-layer $G$-block (the $H$-block agrees to $\pm0.01$). P1 is the
readout cosine weighted by $\|r_j\|\|A_jB_j\|$; nulls pair neuron $j$ with
another neuron's quantities.}
\label{tab:main}
\end{table}

\paragraph{R1: the mechanism is confirmed for $k=3$.} In all single-layer
$k=3$ networks the implied bilinear map is the matrix-multiplication tensor
(P2, cosine $1.000$), per-neuron readouts are parallel to $A_jB_j$
(P1, $0.69$--$0.96$ vs $\approx0$), and neurons are composable (P3, $0.75$--$0.91$
vs $0.49$--$0.53$). The pure quadratic network trains best ($0.14^\circ$).

\paragraph{R2: training concentrates the first layer on $l=1$,} from $28\%$ of
the preactivation variance at initialisation to $80$--$88\%$. This is the irrep that carries
the target, so it is expected here; it is not by itself evidence for irrep
specialisation (see Section~\ref{sec:planned}).

\paragraph{R3: higher irreps cancel extraneous terms.} Removing the $l\ge2$
components of the input increases the error ($0.14^\circ\to2.6^\circ$ for $z^2$),
yet leaves the \emph{bilinear} part of the output unchanged or better
(relative MSE to the bilinear part of $GH$: $4.2\cdot10^{-6}\to1.1\cdot10^{-6}$
for $z^2$; same pattern for SiLU). Hence $l=2,3$ are not used to
compute $GH$; they cancel the pure-$G$/pure-$H$ terms created by squaring the
$l=1$ preactivations ($\ip{A}{G}^2$ has $l=0,1,2$ parts).

\paragraph{R4: cancellation is not by sign-flip pairs.} Only $0$--$27\%$ of
neurons have a partner with $(A,-B)$ (cosine $>0.9$), against $1$--$9\%$ under the
null. The polarisation construction of Section~\ref{sec:mechanism} and the
Waring decomposition used in \cite{marchetti2026sequential} are not what these
networks learn (in this AdamW, standard-initialisation regime).

\paragraph{R5: $k=2$ sits at the capacity bound.} Quadratic networks reach
relative MSE $0.887$ on the bilinear part, vs the bound $8/9$
(Proposition~\ref{prop:capacity}); error $\approx109^\circ$. Going beyond the
bound requires higher-order terms that combine the $l=2$ and $l=1$ features of $G$
(e.g.\ $GSG^\top\cdot Gc=GSc$), which a single quadratic layer does not have.

\paragraph{Caveats.} Few seeds; 1001 epochs ($k=2$ runs are likely not
converged); AdamW rather than SGD; P1--P2 use a second-order Taylor expansion
for SiLU (exact for $z^2$).
Earlier $\SO$ runs in this project used a non-Haar training sampler and a
network with two hidden layers despite \texttt{num\_layers=1}; all numbers here
use the corrected sampler and a true single hidden layer unless stated.

%======================================================================
\section{Selection: the truncated delta as input and target}
\label{sec:planned}
%======================================================================

The matrix target $GH$ contains only $l=1$, so it cannot reveal irrep
\emph{specialisation}: every useful neuron is an $l=1$ neuron. The finite-group
phenomenon concerns targets containing all irreps. Its $\SO$ counterpart is:

\paragraph{Design.}
\begin{itemize}
\item Input $(\delta^L(G),\delta^L(H))$, target $\delta^L(GH)$, with
  $D^l(GH)=D^l(G)D^l(H)$; $L=3$ (width 2048) and $L=4$ (width 4096).
\item One hidden layer with biases, $\sigma(z)=z^2$, no normalisation; MSE loss;
  $5\cdot10^4$ Haar triples ($80/10/10$), batch 128, 1500 epochs, constant
  learning rate.
\item Plain SGD with small ($\times0.1$) or standard initialisation, weight decay
  $\{10^{-5},10^{-4},10^{-3}\}$; AdamW (lr $3\cdot10^{-4}$) for comparison;
  11 runs, two seeds for the main configurations.
\end{itemize}
By Proposition~\ref{prop:capacity} the target is exactly representable, and by
the $\sqrt{d_l}$ weights every irrep block has the same target energy per
dimension as in the true delta.

\paragraph{Normalisation.} Two scale effects must be controlled when the number
of output coordinates changes. (i) The MSE averages over the output coordinates, so
the gradient reaching each irrep block shrinks as $1/N_L$ while weight decay does
not: with $L\ge5$ and the hyperparameters that work for $L\le4$, training never
leaves the initial plateau (training loss exactly $1.0$ for $1500$ epochs). Scaling
the learning rate alone does not help, because SGD's coupled decay
$\eta(\nabla\mathcal L+\lambda\theta)$ keeps the signal-to-decay ratio fixed; scaling
the loss by $s=N_L/N_{L_0}$ (equivalently $\eta\to s\eta$, $\lambda\to\lambda/s$)
restores training. (ii) In the pixel basis, default initialisation and MSE
averaging over $N$ outputs make the effective Fourier-coordinate dynamics
$\sqrt{N/N_L}$ and $N/N_L$ slower; we correct both (Section~R8b).

\paragraph{Per-neuron irrep spectrum (fit-free).} Split neuron $j$'s
$G$-weights, $H$-weights and readout into blocks $w_j^{(l)},v_j^{(l)},r_j^{(l)}$.
With the $\sqrt{d_l}$ input weights, the Haar variance of the $l$-part of its
$G$-preactivation is exactly $\|w_j^{(l)}\|_F^2$ (Section~\ref{sec:encodings}).
Define $p^G_j(l)\propto\|w_j^{(l)}\|^2$, and likewise $p^H_j$, $p^R_j$.
\begin{enumerate}
\item \textbf{Purity} $\max_lp_j(l)$ (and participation ratio $1/\sum_lp_j(l)^2$);
  fraction of neurons above $0.9$.
\item \textbf{Agreement}: the $3$-way confusion matrix of
  $\arg\max p^G_j$, $\arg\max p^H_j$, $\arg\max p^R_j$ (prediction: diagonal).
\item \textbf{Census}: neurons per $l$, vs target energy and $d_l$.
\item \textbf{Nulls}: initialisation; random re-pairing of $(w_j,v_k,r_m)$.
\item \textbf{Causal test}: remove all neurons labelled $l$; only the
  $l$-component of the output error should increase.
\item \textbf{Within an irrep}: with $A_j=w_j^{(l)}$, $B_j=v_j^{(l)}$ (as
  $d_l\times d_l$ matrices), test $r_j^{(l)}\parallel A_jB_j$ and
  $\sum_jr_j^{(l)}\otimes A_j\otimes B_j\propto M^{(l)}$ (Lemma~\ref{lem:cross}).
\item \textbf{Dynamics}: track $p_j(l)$ during training; test whether irreps
  are acquired sequentially and in the order predicted by
  \cite{marchetti2026sequential} (for a one-hot encoding all
  $\|\hat x(\rho)\|_{\mathrm{op}}$ are equal, so the order is by increasing
  dimension, here increasing $l$).
\end{enumerate}

\begin{hypothesis}\label{hyp:pure}
With weight decay, neurons are irrep-pure ($p_j$ concentrated on one $l$) with
the same $l$ in $G$, $H$ and the readout. Rationale: the target is
block-diagonal by irrep, so a mixed neuron produces cross terms in sectors absent
from the target that other neurons must cancel, at extra norm cost.
\end{hypothesis}

\begin{hypothesis}\label{hyp:order}
Under small initialisation, irreps are learned in increasing $l$.
\end{hypothesis}

\subsection{Results}

\begin{table}[h]
\centering\small
\resizebox{\textwidth}{!}{%
\begin{tabular}{lcccccc}
\hline
run & err ($^\circ$) & pure & agree (null) & active neurons per $l$ & $\cos(T_l,M_{d_l})$ & $s_2/s_1(A_j)$\\
\hline
$L{=}4$, SGD, small init, wd $10^{-4}$ & 2.53 & 1.00 & 1.00 (0.21) & 27/125/343/683 & 1.000/1.000/0.994/0.964 & .001/.003/.004/.006\\
$L{=}3$, SGD, std init, wd $10^{-3}$ & 0.53 & 1.00 & 1.00 (0.33) & \textbf{27/125}/323 & 1.000/1.000/0.962 & .001/.003/.005\\
$L{=}3$, SGD, small init, wd $10^{-4}$, s42 & 0.52 & 1.00 & 1.00 (0.33) & 30/130/334 & 1.000/1.000/0.970 & .001/.002/.003\\
$L{=}3$, SGD, small init, wd $10^{-4}$, s43 & 2.56 & 1.00 & 1.00 (0.45) & 27/130/351 & 1.000/1.000/0.976 & .002/.003/.004\\
$L{=}3$, SGD, std init, wd $10^{-4}$ & 0.20 & 1.00 & 1.00 (0.60) & 28/170/1115 & 1.000/0.999/0.995 & .001/.012/.101\\
$L{=}3$, SGD, small init, wd $10^{-5}$ & 0.19 & 1.00 & 1.00 (0.45) & 38/312/1183 & 1.000/0.999/0.996 & .002/.093/.189\\
$L{=}3$, AdamW, wd $10^{-4}$ (2 seeds) & 0.51--0.55 & 1.00 & 1.00 (0.50) & $\approx$43/267/1263 & 1.000/0.997/0.996 & .002/.045/.118\\
$L{=}3$, SGD, small init, wd $10^{-3}$ & 126 & \multicolumn{5}{c}{did not leave the initial plateau}\\
\hline
\end{tabular}}
\caption{Truncated delta, $\sigma=z^2$, after 1500 epochs (converged
counts after 21{,}500 epochs are in R8). ``pure'': fraction of active
neurons (smallest set covering $95\%$ of $\sum_j\|w_j\|\|v_j\|\|r_j\|$) with more
than $90\%$ of their energy in one $l$, minimum over $G$-weights, $H$-weights and
readout ($0.00$ at initialisation). ``agree'': same label in all three; null:
random re-pairing of neurons' $(w,v,r)$. $T_l=\sum_jr_j^{(l)}\otimes A_j\otimes
B_j$ over $l$-labelled neurons; $M_d$ is the $d\times d$ matrix-multiplication
tensor. Rows omitted: SGD lr 1.0 and AdamW wd $10^{-3}$ (same pattern).}
\label{tab:onehot}
\end{table}

\paragraph{R6: irrep specialisation (Hypothesis~\ref{hyp:pure}) holds in every
run that learned} (10/11, SGD and AdamW, both initialisations, all weight decays):
all active neurons are pure, with the same $l$ in $G$, $H$ and readout (diagonal
confusion matrices).

\begin{figure}[h]
\centering
\includegraphics[width=\textwidth]{figures/purity.pdf}
\caption{Selection of irreps ($L=4$, SGD, $21{,}500$ epochs). (a) At initialisation,
each neuron's first-layer $G$-weights spread their energy over the irrep blocks in
proportion to the block sizes $d_l^2$. (b) After training, every active neuron has
all its energy in a single irrep. (c) Number of active neurons per irrep (sharp
gap in contribution), against $d_l^3$ and $d_l^3+d_l$.}
\label{fig:purity}
\end{figure}

\begin{figure}[h]
\centering
\includegraphics[width=0.9\textwidth]{figures/harmonics.pdf}
\caption{Spherical harmonics read out from the network ($L=4$ run). For each $l$, a
linear decoder from the preactivations of the $l$-labelled neurons alone to
$D^l(G)$ (held-out $R^2\ge0.99998$) is evaluated at rotations sending the north pole
to each point of the sphere; its zonal column is $\sqrt{4\pi/d_l}\,Y_l^m$ (upper
rows), shown above the exact harmonics (lower rows) with a shared colour scale per
$(l,m)$. Since the first layer is linear in $\delta^L(G)$, this shows availability
(the $l$-labelled neurons span the $l$-block), not selection; selection is
Figure~\ref{fig:purity}.}
\label{fig:harmonics}
\end{figure}

\paragraph{R7: each irrep block is a matrix-multiplication circuit:}
$\cos(T_l,M_{d_l})=0.96$--$1.000$ and $r_j^{(l)}\parallel A_jB_j$ (weighted
cosine $0.94$--$1.00$). Ablating the $l$-labelled neurons destroys only the
$l$-component (e.g.\ $L=4$, removing $l=2$: relative MSE $0.01/4.58/0.01/0.03$).

\paragraph{R8: under weight decay, $\approx d_l^3$ rank-one product units.}
After 21{,}500 epochs of SGD with weight decay, each irrep shows a sharp gap in
neuron contribution (the next neuron is at $\sim10^{-119}$ of the largest). The
number of units is $27=3^3$ and $125=5^3$ for $l=1,2$, and $350=7^3+7$,
$738=9^3+9$ for $l=3,4$; the count $350$ is reproduced by two runs with different
initialisation, learning rate and weight decay. (At epoch 1500 a coverage
threshold gave $343$ for $l=3$, an artefact of the threshold.) We do not yet
understand the extra $d_l$ units for $l\ge3$. Each surviving neuron computes
$(a^\top D^l(G)u)(u^\top D^l(H)b)$ and writes it along $ab^\top$: $A_j$, $B_j$,
$r_j$ are rank one ($s_2/s_1\le0.006$), $A_j$'s right and $B_j$'s left singular
vectors coincide ($|\cos|=1.000$), and all units of an irrep have the same norm
(spread $\le0.4\%$). This is not a rotated copy of~\eqref{eq:matmul}: for $l=1$
the units split into 3 nearly orthogonal output directions $a$ with 9 units each,
but the inner vectors take 17 directions. Each output row is computed by its own
9-term vector-by-matrix product, which is rank-minimal. AdamW, whose decoupled decay is negligible
here, keeps $\approx2$--$4\times d_l^3$ less clean units.

\paragraph{R8e: the $d_l^3$ count is the minimum-norm decomposition.} Weight decay on
a one-hidden-layer $z^2$ network penalises $\sum_j(\|r_j\|^2+\|a_j\|^2+\|b_j\|^2)$, and
the bilinear part of its output in block $l$ is $\sum_j r_j\otimes a_j\otimes b_j$,
which must equal the $d_l\times d_l$ matrix-multiplication tensor $M_{d_l}$. Solving
directly, without a network,
\[
\min_{(r_j,a_j,b_j)_{j\le3d^3}}\ \Bigl\|\sum_j r_j\otimes a_j\otimes b_j-M_d\Bigr\|^2
+\lambda\sum_j\bigl(\|r_j\|^2+\|a_j\|^2+\|b_j\|^2\bigr),
\]
from $3d^3$ random units ($\lambda=3\cdot10^{-3}$, 3 seeds), leaves exactly
$27,125,343,729$ units for $d=3,5,7,9$, with equal contributions (spread
$\le0.3\%$) and every other unit at contribution $\le10^{-26}$. The count of the
trained networks is therefore a property of the weight-decay objective on the target
block, not of the training dynamics, and the minimum-norm decomposition is not a
low-rank (Strassen-type) one: the $3\times3$ tensor has rank at most $23$
\cite{laderman1976}. This computation ignores the non-bilinear terms that a $z^2$
network must also cancel; the extra $d_l$ units observed for $l\ge3$ are a natural
candidate for that role. Script: \texttt{scripts/min\_norm\_decomposition.py}.

\paragraph{R8b: hidden irreps (pixel basis).} We repeat the experiment with
input and target encoded by the truncated delta sampled at $N$ Haar-random anchors,
$x_G[i]=K_L(G,R_i)/\sqrt N$ with the $l=0$ term included
(Section~\ref{sec:computing-encoding}); no irrep is visible in these coordinates.
Two scale corrections make the runs comparable with the Fourier-basis ones: the
default initialisation is rescaled by $\sqrt{N/(N_L+1)}$, since $1/\sqrt{\text{fan-in}}$
would otherwise make the effective Fourier weights $W\Phi$ that much smaller, and
the MSE is reported per Fourier coordinate (multiplied by $N/N_L$), since averaging
over $N$ outputs would shrink gradients by that factor. Analysis maps the weights
back to Fourier blocks ($W_G\Phi$, $W_H\Phi$, $\Phi^+W_{\mathrm{out}}$).

\begin{center}\small
\begin{tabular}{llllll}
\hline
$L$ & $N$ & err ($^\circ$) & pure / agree & active neurons per $l$ & $\cos(T_l,M_{d_l})$\\
\hline
2 & 256--1024 & 0.63--0.33 & 1.00 / 1.00 & 27--30 / 122 & 1.000 / 0.97\\
2 & 2048 (lr 0.05) & 0.22 & 1.00 / 1.00 & 27 / 127 & 1.00 / 0.98\\
3 & 256--2048 & 3.83--0.84 & 1.00 / 1.00 & 27 / 125--130 / 322--325 & 1.000 / 1.000 / 0.96\\
3 & 4096 (lr 0.05) & 0.56 & 0.98 / 1.00 & 30 / 130 / 316 & 1.00 / 1.00 / 0.96\\
4 & 1024--2048 ($L{=}4$ config.) & 2.23--2.29 & 1.00 / 1.00 & 27--30 / 125 / 343--350 / 697--733 & 1.00 / 1.00 / 1.00 / 0.97\\
3, Fourier basis & --- & 0.53 & 1.00 / 1.00 & 27 / 125 / 323 & 1.000 / 1.000 / 0.962\\
\hline
\end{tabular}
\end{center}

Whenever training succeeds, the structure is the same as in the Fourier basis:
all active neurons are irrep-pure and consistent, the neuron counts and the
matrix-multiplication alignment match the Fourier-basis run with the same
hyperparameters, and weight decay removes all first-layer weight outside
$\operatorname{range}(\Phi)$ (100\% of the norm ends up in the $(N_L+1)$-dimensional
subspace occupied by the data). The number of anchors only affects accuracy,
through the Riemann sum. This is the $\SO$ counterpart of the finite-group
situation: the irreps are not visible in the input, yet the hidden units select
them. At learning rate $0.2$, the runs with the largest corrections ($N\ge2048$
for $L=2$, $N=4096$ for $L=3$) diverge; at $0.05$ (or $0.02$) they train and show
the same structure (see R8c for seeds). For $L=4$ the hyperparameters of the $L=3$
runs leave training on the initial plateau, as in the Fourier basis; with those that
work for $L=4$ in the Fourier basis (small initialisation, lr $0.5$, weight decay
$10^{-4}$), the pixel-basis runs reproduce the Fourier-basis structure, including
$343$--$350$ units for $l=3$.

\paragraph{R8c: robustness across seeds.} Each configuration below was trained
with three seeds for $5000$ epochs (for the pixel basis the anchors change with the
seed). All runs have $100\%$ irrep-pure active neurons with consistent labels.

\begin{center}\small
\begin{tabular}{lllll}
\hline
configuration & err ($^\circ$) & pure / agree & active neurons per $l$ & min $\cos(T_l,M_{d_l})$\\
\hline
$L{=}3$, Fourier, standard init & 0.52--0.64 & 1.00 / 1.00 & 27 / 125--130 / 322--329 & 1.00 / 1.00 / 0.96\\
$L{=}3$, Fourier, small init & 0.65--2.37 & 1.00 / 1.00 & 27--30 / 125--130 / 326--332 & 1.00 / 0.98 / 0.96\\
$L{=}4$, Fourier, small init & 1.42--4.00 & 1.00 / 1.00 & 27 / 125 / 343--350 / 668--676 & 1.00 / 1.00 / 0.99 / 0.95\\
$L{=}3$, pixel, $N=2048$ & 0.80--0.90 & 1.00 / 1.00 & 27 / 130 / 322--323 & 1.00 / 1.00 / 0.96\\
$L{=}2$, pixel, $N=1024$ & 0.30--0.33 & 1.00 / 1.00 & 27 / 122 & 1.00 / 0.97\\
\hline
\end{tabular}
\end{center}

The counts for the lower irreps are reproduced exactly ($27$ in every run; $125$
or $130$ for $l=2$); the highest irrep of each run, which carries the largest target
block and converges last, is below its long-run value after $5000$ epochs
(compare R8).

\paragraph{R8d: larger band-limits, and are all irreps learned?} With width at
least $\sum_{l\le L}d_l^3$ and the loss rescaled as described under Normalisation
(training otherwise stays on the initial plateau for $L\ge5$), irreps are acquired
in increasing order of $l$, and every irrep up to $L$ is eventually learned when
training converges. Since the geodesic error only measures the $l=1$ block, we
report the relative error of each block. Every block is fitted (relative MSE
$<0.1$) in all converged runs for $L\le4$. For $L=5$ after $4000$ epochs, blocks
$l\le4$ are fitted in all three runs and $l=5$ is fitted in one and at $0.16$--$0.19$
in the other two; neurons are $100\%$ pure with counts
$27/125/341$--$350/724$--$729$ ($9^3=729$ for $l=4$). For $L=6$, after $1500$ epochs
only $l=1$ (or $l\le2$) is fitted; after $4000$ epochs two of three runs have fitted
$l\le4$ (relative MSE $\le0.14$) with $l=5,6$ at $0.37$--$0.49$, and the third has
not progressed beyond $l=1$. In these runs the readout weights are already pure
($0.99$--$1.00$) and input purity rises with training (median $0.89$--$0.90$).
Higher irreps are thus learned later, not skipped; within our budget the largest
ones were not reached for $L\ge6$, and $L=8,10$ did not train.

\paragraph{R9: the regime matters.} Small initialisation with weight decay
$10^{-3}$ never leaves the initial plateau. Without strong weight decay, the census
keeps shrinking long after the error is small (main run, $l=3$: $955\to334$
between epochs 240 and 1500). Continuing two runs to 21{,}500 epochs leaves the
census and the rank-one structure unchanged after $\sim5000$ epochs but sharpens
the patterns: unit norms within an irrep equalise ($L=4$, $l=2$: spread
$0.83\to0.006$), and the three output directions of the $l=1$ units become
orthogonal (max $|\cos|$ $0.15\to0.01$).

\paragraph{R9b: explicit weight decay is needed for pruning.} Without weight decay
($L=3$, SGD) the network is still irrep-pure: $100\%$ of active neurons are pure
from initialisation $\times0.1$, and $80$--$90\%$ (median purity $0.95$) from
standard initialisation, with consistent labels in $99$--$100\%$ of neurons. But it
keeps $1544$--$1840$ active neurons, against $\approx\sum_ld_l^3=475$ with weight
decay. This does not show that purity is independent of norm control: a small
initialisation is itself an implicit regularisation, since gradient flow from near
zero is biased towards low-norm solutions, much like the limit of a vanishing
Lagrange multiplier on the norm constraint that weight decay imposes. What these runs
show is that implicit norm control suffices for purity, while explicit weight decay
is needed to reach the minimum-norm count (R8e). Separating purity from norm control
requires the lazy regime (large initialisation, no weight decay), where this
implicit bias disappears. There, with initialisation $\times3$ and $\times10$, the
network partly fits the target (relative error $0.29$--$0.35$ and $0.82$--$0.84$
per block after $1500$ epochs) while its neurons keep their initial mixing (no pure
neuron; median purity $0.59$, the value at initialisation). Without any norm
control, explicit or implicit, the hidden units are not selected; these lazy runs
are not converged, and larger learning rates diverged.

\paragraph{R10: learning order.} With per-neuron scaled SGD
($\eta_i\propto\|\theta_i\|^{-1}$) from a vanishing initialisation ($\times10^{-3}$ or
$\times10^{-4}$) and no weight decay, the time-rescaled gradient flow of
\cite{marchetti2026sequential,kunin2025agf}, the irrep blocks are learned one after
another, in increasing $l$: the relative error of block $l$ falls below $0.5$ at
epochs $26/31/35$ for $L=3$ and $59/71/80/89$ (initialisation $10^{-4}$) or
$50/60/67/73$ ($10^{-3}$) for $L=4$. This is the staircase predicted by AGF for a
one-hot-like encoding, ordered by irrep dimension (Hypothesis~\ref{hyp:order}).
At the end of training the neurons are pure (median purity $0.97$--$0.99$) with
consistent labels in all neurons. With standard SGD and initialisation $\times0.1$
the same order appears, but the steps overlap.

\paragraph{Availability versus selection.} The one-hot encoding of a finite group
is the regular representation written in the position basis, and a first layer
with isotropic initialisation and weight decay cannot distinguish it from the
Fourier basis (Section~\ref{sec:finite}). Irrep matrix elements are therefore
available from the start; a network does not have to ``discover'' the Fourier
transform. The pixel-basis experiment (R8b) makes this concrete for $\SO$: whether
the truncated delta is presented through its Fourier coefficients or through
samples at random anchors, the trained network is the same up to the change of
basis. What is not given by the encoding is the \emph{choice of hidden basis}.
With an elementwise activation, rotating the hidden units changes the function, so
some basis must be chosen; at initialisation every unit mixes all irreps
(Figure~\ref{fig:purity}a). That every surviving unit ends up in a single irrep,
consistently in its two input blocks and its readout, and that each block is then
served by equal-norm rank-one product units, is the non-trivial outcome. Our
experiments distinguish two effects of norm control. Implicit control (small
initialisation) suffices for purity, while explicit weight decay is needed to prune
to the minimum-norm count (R9b); whether purity can arise without any norm control
is tested by the lazy regime. We propose the following mechanism for purity in the
rich regime: by Schur orthogonality the
gradient of a unit decouples across irreps at leading order, each irrep component
of a unit grows as a degree-3 homogeneous system (rich get richer), so the
best-aligned irrep runs away first, and cross terms between irreps, which only add
error, are suppressed. This is the utility-maximisation step of AGF
\cite{kunin2025agf,marchetti2026sequential}, and is consistent with the ordered
acquisition we observe (R10) and with purity being weaker from larger
initialisations. Weight decay then prunes the pure units to the minimum-norm
decomposition of each block, $d_l^3$ units (R8e). Making the purity argument
precise is left open.

\paragraph{Relation to Marchetti et al.} Our setting is their two-layer,
quadratic, pairwise composition \cite{marchetti2026sequential}, transported to
$\SO$. Two of their ingredients appear in trained networks: the sigma-pi-sigma
product term, and the need to cancel extraneous terms. Their explicit Waring
construction, with sign-flip pairs of units, is not what our networks learn (R4).
Their theorem assumes that every Fourier coefficient $\hat x(\rho)$ of the encoding
is invertible or zero, which always holds for generic encodings in $\R[G]$;
Proposition~\ref{prop:capacity} describes the rank-deficient case that can occur
for $\SO$ encodings. Their construction uses rank-one matrices, which they relate
to the $\rho$-sets of Wu et al.\ \cite{wu2024unifying}; the units we find after
training are rank one in the same sense.

\paragraph{Ordering and smoothness.} The $D^l$ are eigenfunctions of the
Laplace--Beltrami operator of $\SO$ with eigenvalue $l(l+1)$, and the degree
filtration (Proposition~\ref{prop:degree}) gives the same order. A preference for
low $l$ could therefore come from a smoothness (Dirichlet-energy) bias, from the
dimension ordering of AGF \cite{marchetti2026sequential}, or from polynomial
degree when irreps must be computed from rotation matrices. These predictions
coincide for the truncated delta; reweighting the target blocks would separate
them, since the AGF order follows the weights while a smoothness bias does not.

\paragraph{Choice of activation.} We use $\sigma(z)=z^2$ because it makes the
network a polynomial whose cross terms are exactly the products
$\ip{A_j}{D^l(G)}\ip{B_j}{D^l(H)}$, so that every claim above can be checked in
closed form. Other activations lead to different, approximate, circuits on the
same tasks \cite{zhong2023clock, nanda2023progress}; we expect the separation of
availability and selection to carry over, but do not study it here.

\paragraph{Limitations.} All networks have one hidden layer and are trained with
MSE. The structural results hold in the regime where training succeeds; small
initialisation with strong weight decay stays on the initial plateau, and the
learning rate, initialisation and loss scale must be adapted to the encoding
(Section~\ref{sec:planned}, Normalisation). The neuron counts are observed, not
derived, and the extra $d_l$ units for $l\ge3$ are unexplained. The sharp
truncation of the delta converges only in the sense of distributions
(Figure~\ref{fig:delta}b).

\paragraph{Open questions.} How the extraneous terms are cancelled at the
population level (R3, R4); why weight decay selects $\approx d_l^3$ units per irrep
rather than a lower-rank decomposition (the $3\times3$ matrix-multiplication tensor
has rank at most $23$ \cite{laderman1976}, yet R8 finds 27), and what the extra $d_l$
units for $l\ge3$ do; and how the picture changes with depth, in particular when
irreps must be computed from rotation matrices (Proposition~\ref{prop:degree}).

%======================================================================
\section*{Reproducibility}
%======================================================================
{\sloppy
All code is in the NeuralRep repository, with tests. Encodings:
\texttt{neuralrep/datamodules/so3\_onehot.py} (truncated delta in Fourier and
pixel bases), \texttt{neuralrep/algorithms/callbacks/wigner\_d.py} (real Wigner
matrices). Training: \texttt{scripts/train\_so3\_onehot.py} (truncated delta),
\texttt{scripts/train\_so3\_mechanism.py} (matrix target $GH$); cluster job
scripts in \texttt{scripts/slurm/}, with the exact configuration lists used.
Analysis: \texttt{scripts/analyze\_so3\_onehot.py} and
\texttt{neuralrep/algorithms/callbacks/so3\_irrep\_spectrum.py} (fit-free irrep
spectra, purity, counts, matrix-multiplication alignment, product units),
\texttt{scripts/analyze\_so3\_mechanism.py} and
\texttt{neuralrep/algorithms/callbacks/so3\_mechanism.py} (Casimir projectors, $l=1$
maps, output sectors). Availability bound: \texttt{scripts/capacity\_check.py}.
Figures: \texttt{scripts/plot\_delta\_kernel.py}, \texttt{scripts/plot\_so3\_purity.py},
\texttt{scripts/plot\_so3\_harmonics.py}. Per-run configurations and analyses are
in \texttt{results/}, documented in \texttt{SO3-RESULTS.md}.\par}

\begin{thebibliography}{99}
\bibitem{marchetti2026sequential} G.~L.~Marchetti, D.~Kunin, A.~Myers, F.~Acosta,
  N.~Miolane. Sequential group composition: a window into the mechanics of deep
  learning. arXiv:2602.03655, 2026.
\bibitem{kunin2025agf} D.~Kunin, G.~L.~Marchetti, F.~Chen, D.~Karkada, J.~B.~Simon,
  M.~R.~DeWeese, S.~Ganguli, N.~Miolane. Alternating gradient flows: a theory of
  feature learning in two-layer neural networks. arXiv:2506.06489, 2025.
\bibitem{chughtai2023toy} B.~Chughtai, L.~Chan, N.~Nanda. A toy model of
  universality: reverse engineering how networks learn group operations. ICML 2023. arXiv:2302.03025.
\bibitem{nanda2023progress} N.~Nanda, L.~Chan, T.~Lieberum, J.~Smith,
  J.~Steinhardt. Progress measures for grokking via mechanistic
  interpretability. ICLR 2023. arXiv:2301.05217.
\bibitem{gromov2023grokking} A.~Gromov. Grokking modular arithmetic. arXiv:2301.02679, 2023.
\bibitem{zhong2023clock} Z.~Zhong, Z.~Liu, M.~Tegmark, J.~Andreas. The clock
  and the pizza: two stories in mechanistic explanation of neural networks.
  NeurIPS 2023. arXiv:2306.17844.
\bibitem{stander2024cosets} D.~Stander, Q.~Yu, H.~Fan, S.~Biderman. Grokking
  group multiplication with cosets. arXiv:2312.06581, 2024.
\bibitem{morwani2023feature} D.~Morwani, B.~L.~Edelman, C.-A.~Oncescu, R.~Zhao,
  S.~Kakade. Feature emergence via margin maximization: case studies in
  algebraic tasks. arXiv:2311.07568, 2023.
\bibitem{wu2024unifying} W.~Wu, L.~Jaburi, J.~Drori, J.~Gross. Unifying and verifying
  mechanistic interpretations: a case study with group operations.
  arXiv:2410.07476, 2024.
\bibitem{laderman1976} J.~D.~Laderman. A noncommutative algorithm for
  multiplying $3\times3$ matrices using 23 multiplications. Bull.\ AMS, 1976.
\end{thebibliography}

\appendix
\section{Real Wigner matrices}\label{app:wigner}

\paragraph{Definition.} Let $\mathcal Y_l$ be the space of spherical harmonics of
degree $l$ (restrictions to $S^2$ of harmonic homogeneous polynomials of degree $l$
on $\R^3$), $\dim\mathcal Y_l=2l+1$. $\SO$ acts on $\mathcal Y_l$ by
$(G\cdot f)(x)=f(G^{-1}x)$; this action is irreducible, and every irrep of $\SO$ is
of this form for exactly one $l\ge0$. With the real orthonormal basis
$(Y_l^m)_{m=-l}^{l}$ of $L^2(S^2)$,
\[
Y_l^0=N_{l0}P_l(\cos\theta),\quad
Y_l^m=\sqrt2N_{lm}P_l^{m}(\cos\theta)\cos(m\varphi),\quad
Y_l^{-m}=\sqrt2N_{lm}P_l^{m}(\cos\theta)\sin(m\varphi)\quad(m>0),
\]
with $P_l^m$ the associated Legendre functions and $N_{lm}$ the usual normalisation,
the real Wigner matrix is defined by
$Y_l^m(G^{-1}x)=\sum_{m'}D^l_{m'm}(G)Y_l^{m'}(x)$. It is orthogonal, and
$D^l(GH)=D^l(G)D^l(H)$.

\paragraph{Euler angles.} Writing $G=R_z(\alpha)R_y(\beta)R_z(\gamma)$, the complex
Wigner matrix in the basis of complex harmonics is
$\mathcal D^l_{mm'}(\alpha,\beta,\gamma)=e^{-im\alpha}d^l_{mm'}(\beta)e^{-im'\gamma}$,
with $d^l$ the Wigner small-$d$ matrix; the real Wigner matrix is obtained from it by
the fixed unitary change of basis between complex and real harmonics.

\paragraph{Properties used in the paper.} (i) Schur orthogonality,
equation~\eqref{eq:schur}. (ii) The zonal column is a vector of harmonics:
$D^l(G)_{\cdot,0}=\sqrt{4\pi/d_l}\,(Y_l^m(G\hat z))_m$; this is what
Figure~\ref{fig:harmonics} displays. (iii) The character is
$\tr D^l(G)=\sin\bigl((2l+1)\theta/2\bigr)/\sin(\theta/2)$, with $\theta$ the rotation
angle of $G$. (iv) The entries of $D^l(G)$ are polynomials of degree exactly $l$ in
the entries of $G$ (Proposition~\ref{prop:degree}).

\paragraph{Computation.} We compute $D^l(G)$ from the definition: evaluate the real
harmonics on a quadrature grid of $S^2$ fine enough for degree $l$, evaluate them at
the rotated grid points $G^{-1}x$, and solve the linear system for the coefficients
$D^l_{m'm}(G)$. We checked that the result is an orthogonal homomorphism to
machine precision ($\max|D^l(GH)-D^l(G)D^l(H)|\le4\cdot10^{-13}$ and
$\E[D^l_{ij}(G)^2]\,d_l=1$ for $l\le16$), and that its zonal column matches
$\sqrt{4\pi/d_l}\,Y_l^m(G\hat z)$ to $2\cdot10^{-14}$ for $l\le3$
(\texttt{neuralrep/algorithms/callbacks/wigner\_d.py}).

\end{document}
